% =============================================================================
% TEMPLATE TESI - Politecnico di Torino
% =============================================================================

\documentclass[11pt,a4paper]{report}
\usepackage[top=3cm, bottom=2.5cm]{geometry}

% --- Encoding & Language ---
\usepackage[utf8]{inputenc}
\usepackage[T1]{fontenc}
\usepackage[english]{babel}         % Cambiare in [italian] se la tesi è in italiano

% --- Typography ---
\usepackage{microtype}

% --- Hyperlinks ---
\usepackage{hyperref}
% \hypersetup{
%     pdfpagemode={UseOutlines},
%     bookmarksopen,
%     pdfstartview={FitH},
%     colorlinks,
%     linkcolor={blue},
%     citecolor={blue},
%     urlcolor={blue}
% }

% --- Tables ---
\usepackage{booktabs}
\usepackage{multirow}

% --- Graphics ---
\usepackage{graphicx}
\usepackage{wrapfig}

% --- Code Listings ---
\usepackage{listings}
\usepackage{minted}

% --- Colors ---
\usepackage{color}
\usepackage{xcolor}
\definecolor{lbcolor}{rgb}{0.9,0.9,0.9}

% --- Quotations ---
\usepackage{csquotes}

% --- Math ---
\usepackage{amsmath, amssymb, stmaryrd}

% --- Algorithms ---
\usepackage[ruled,vlined]{algorithm2e}

% --- Utility ---
\setlength{\marginparwidth}{2cm}
\usepackage{todonotes}

% =============================================================================
% Comandi personalizzati - Definire qui le abbreviazioni specifiche della tesi
% =============================================================================
% Esempio:
% \newcommand{\myterm}{\texttt{MyTerm}}

% =============================================================================
% Path immagini
% =============================================================================
\graphicspath{ {./images/} }